\documentclass{report}  
\usepackage[utopia]{mathdesign} 
%\usepackage{amsmath,amsfonts,amsthm,amssymb,mathtools}

\input{preamble.tex}


\usepackage{amsmath,amsthm,mathtools}
\usepackage{titlesec}
\usepackage{microtype}
\definecolor{myblue}{RGB}{0,82,155}
\usepackage{paracol}




%\usepackage[scr]{rsfso}

%\usepackage{libertine}
%\usepackage{mathpazo}
%\usepackage{palatino}
%usepackage{crimson}


\newcommand{\faketarget}{\oplus\!\!\!\!\odot}
\newcommand{\target}{%
  \begin{tikzpicture}[scale=0.5]
    \fill[black] (0,0) circle (0.1);
    \draw (0,0) circle (0.2);
    \draw (0,0) circle (0.3);
  \end{tikzpicture}%
}

\usepackage[clock]{ifsym}


\title{\Huge{Interface Personne Marchine}\\{IFT2905}\\{\textbf{Introduction}}}
\author{\huge{Franz Girardin}}
\date{\today}
\lstset{inputencoding=utf8/latin1}

            %%%%%%%%%%%%%%%%%  Sect.                          %%%%%%%%%%%%%%%%%%%%%%%%%%%%%%%%%%%%%%%%%%%%%%%%%%%%%%%%%
\usepackage{helvet}
\titleformat{\chapter}[display]
  {\normalfont\bfseries\color{myblue}}
  {\filleft%
    \begin{tikzpicture}
    \node[
      outer sep=0pt,
      text width=1.5cm,
      minimum height=2cm,
      fill=myblue,
      font=\color{white}\fontsize{40}{50}\selectfont,
      align=center
      ] (num) {\thechapter};
    \node[
      rotate=90,
      anchor=south,
      font=\color{black}\Large\normalfont
      ] at ([xshift=-5pt]num.west) {\textls[180]{\textsc{Section}}};  
    \end{tikzpicture}%
  }
  {5pt}
  {\titlerule[2.0pt]\vskip3pt\titlerule\vskip4pt\large\normalfont}

\titleformat{\section}
  {\normalfont\scshape}{\thesection}{1em}{}



\titleformat{\section}
  {\normalfont\scshape}{\thesection}{1em}{}


% Customizing the spacing for the chapter titles
\titlespacing*{\chapter}{0pt}{0pt}{20pt}

%\usepackage[utopia]{mathdesign}
% Allow hfill in math environment
\newcommand{\specialcell}[1]{\ifmeasuring@#1\else\omit$\displaystyle#1$\ignorespaces\fi}

% Allow you to do the non implication (implication barred)
\newcommand{\notimplies}{%
  \mathrel{{\ooalign{\hidewidth$\not\phantom{=}$\hidewidth\cr$\implies$}}}}



\DeclareRobustCommand{\looongrightarrow}{%
  \DOTSB\relbar\joinrel\relbar\joinrel\relbar\joinrel\rightarrow
}


\title{\Huge{Interface PM}\\{IFT2905}\\{\textbf{Travaux pratiques}}}
\author{\huge{Franz Girardin}}
\date{\today}


\begin{document}
\maketitle
\pagebreak
\tableofcontents
\pagebreak

\begin{multicols*}{3}
\chapter{TP1 Notation} 

    \footnotesize

\begin{Exercice}{}{}
Le but de cet exercice est de comprendre comment passer d’une notation in-
fixe (lisible par les humains) à une notation postfixe et préfixe. Ces notations
sont utiles car elles permettent à un ordinateur de les ´evaluer (calculer) plus
facilement et sans ambiguité.   
\end{Exercice}

\paragraph{1.1}



\begin{align*}
    a + b + c = ((a \textcolor{myg}{+}   b) \textcolor{myr}{+}   c) 
             &= ((+ a \, b) + c) \\ 
             &= + (+ a \, b) \; c \\ 
             &= + + a \, b \, c \; 
             \\
             & \textbf{préfixé}  
\end{align*}
\begin{align*}
    a + b + c = ((a \textcolor{myg}{+}   b) \textcolor{myr}{+}   c)           
    &= ((a \, b \, +) + c) \\
    &= ((a \, b \, +) \, c +) \\ 
    &=  a \, b \, + c \, + \\
    & \textbf{postfixé}
\end{align*}


\paragraph{1.2}
\begin{align*}
    a + (b + c) = (a + (+ b \, c)) 
             &= ( + a \, (+ \, b \, c) ) \\
             &= + \; a  +  b \, c \\
             & \textbf{préfixé}  
\end{align*}


\begin{align*}
    a + (b + c) = (a + (b \, c \, +)) 
             &= (a \, ( b \, c) \, + \, + ) \\
             &= a  b \, c \, + +\\
             \\
             & \textbf{postfixé}  
\end{align*}


\paragraph{1.2}
\begin{align*}
    a \times b + c \times d &= (a \times b) + (c \times d) 
                    \\
                  &= (\times a \, b) + (\times c \, d) 
                    \\
                  &= + (\times a \, b) \times c \, d
                    \\
                  &= + \times a \, b \times c \, d
                    \\
                  & \textbf{préfixé}
\end{align*}    

\begin{align*}
    a \times b + c \times d &= (a \times b) + (c \times d) 
    \\
            &= (a \, b \times) + (c \, d \times)
    \\
            &= (a \, b \times) (c \, d \times) + 
    \\ 
            &= a \, b \times c \, d \times +
    \\
            & \textbf{postfixé}
\end{align*}


\begin{align*}
    a + b < a \times (c + d) &= (+ a \, b) < \bigl( a \times (+ c \, d) \bigr)      
    \\ 
            &= (+ a \, b ) < \bigl( \times a (+ c \, d)\bigr)
    \\ 
            &= < (+ a \, b )  \bigl( \times a (+ c \, d)\bigr)
    \\
            &= < + a \, b  \times a + c \, d
    \\
                  & \textbf{préfixé}
\end{align*}


\begin{align*}
    a + b < a \times (c + d) &= (a \, b +) < \bigl( a \times (c \, d +) \bigr)
    \\ 
            &= (a \, b +) < \bigl( a  (c \, d +) \times \bigr)
    \\
            &= (a \, b +)  \bigl( a  (c \, d +) \times < \bigr)
    \\
            &= a \, b +  a  c \, d + \times < \bigr)
            \\
            & \textbf{postfixé}
\end{align*}


\begin{align*}
    &(-b + sqrt(b\times b -4 \times a \times c))/(2 \times a)
    \\ 
    &= 
    (-b + sqrt((\times b \, b) - (\times \times 4\,  a \, c)  ))/(\times 2 \, a)
    \\
    &=\bigl(+ -b \; sqrt\bigl( - \times b \, b  \times \times 4\,  a \, c   \bigr)  \bigr)/(\times 2 \, a)
    \\
    &=\bigl(+ -b \; sqrt\bigl( - \times b \, b  \times \times 4\,  a \, c   \bigr)  \bigr) \times 2 \, a /
    \\ 
    &= + -b \; sqrt - \times b \, b  \times \times 4\,  a \, c -  / \times 2 \, a 
    \\
    & \textbf{préfixé}
\end{align*}


\begin{align*}
    &(-b + sqrt(b\times b -4 \times a \times c))/(2 \times a)
    \\ 
    &= 
    (b- + sqrt((b \, b \times) - (4 \,  a \times  c \, \times)  ))/(2 \, a \times)
    \\
    &=(b- + sqrt( b \, b \times  4 \,  a \times  c \, \times -  ))/(2 \, a \times)
    \\
    &=(b- + b \, b \times  4 \,  a \times  c \, \times -  sqrt )  /(2 \, a \times)
    \\
    &=(b- b \, b \times  4 \,  a \times  c \, \times -  sqrt +) + /(2 \, a \times)
    \\
    &= b- b \, b \times  4 \,  a \times  c \, \times -  sqt + 2 \, a \ \times
    \\
            & \textbf{postfixé}
\end{align*}

\begin{Exercice}{}{}
Le but de cet exercice est de comprendre comment un ordinateur réesout un
calcul à partir d’une notation postfixe.   
\end{Exercice}


\begin{align*}
    ab \textcolor{myp}{+}   c \textcolor{red}{+} : \quad & \\  
    \\
                     & \texttt{stack}(a); 
    \\
                     & \texttt{stack}(b); 
    \\
                     & \textbf{do}   :
                     \\ 
                     & \quad \textcolor{myb}{\texttt{\textit{ // Opération a \textcolor{myp}{+}   b}}}  
    \\
                     & \quad \texttt{stack}\bigl(\texttt{unstack}() \textcolor{myp}{+}   \texttt{unstack}()\bigr);
    \\ 
                     & \texttt{stack}(c) 
    \\
                     & \textbf{do} : 
                     \\ 
                      & \quad \textcolor{myb}{\texttt{\textit{ // Opération (a \textcolor{red}{+}   b) + c }}}                    
    \\
                     & \quad \texttt{stack}\bigl(\texttt{unstack}() \textcolor{myp}{+}   \texttt{unstack}() \bigr)
\end{align*}


\begin{Exercice}{}{}
Donner un exemple d'expression ambigue    
\end{Exercice}

\texttt{\textbf{if}   p \textbf{then} \textbf{if} q \textbf{then} r \textbf{else} s}  

L'expression est ambvigue puisqu'on ne sait pas si le bloc \texttt{else} est associé à la condition 
1, \texttt{if p}, ou à la condition 2 \texttt{if q}.      


\begin{Exercice}{}{}
    Donner une grammaire non ambigue qui associe les \texttt{else} avec les 
    \texttt{if} le plus proche, cmme le font les langage de programmation habituels.  
\end{Exercice}


% Première Grammaire
\begin{align*}
 S ::= X \\
    \mid \text{if} \; E \; \text{then} \; S \; \text{end} \\
    \mid \text{if} \; E \; \text{then} \; S \; \text{else} \; S \; \text{end}           
\end{align*}

Cette grammaire force le \texttt{else} à être explicitement inclus avec son \texttt{if} 
correspondant. Chaque bloc est finalisé par la déclaration \texttt{end}, assurant ansi qu'un 
else peut uniquement être lié au \texttt{if}   le plus proche.

% Deuxième Grammaire
\begin{align*}
 S ::= X \\
    \mid \text{if} \; ( \; E \; ) \; \{ \; S \; \} \\
    \mid \text{if} \; ( \; E \; ) \; \{ \; S \; \} \; \text{else} \; \{ \; S \; \}
\end{align*}

Cette grammaire utilise des parenthèses et des accolades pour structurer les blocs, ce qui rend 
l'association entre if et else explicite et sans ambiguïté.

\chapter{EBNF}

\begin{Exercice}{}{}
 Voici un exemple de programme dans un langage similaire `a Javascript et 
 la grammaire du langage en notation EBNF 
 Compl´eter la grammaire pour qu’elle accepte le programme. 
\end{Exercice}

\begin{adjustbox}{width=\linewidth}
  \begin{lstlisting}
var ackermann = function (m, n) {
  if (m == 0) {
    return n + 1;
  } else if (n == 0) {
      return ackermann (m - 1, 1);
  } else {
      return ackermann (m - 1, ackermann (m , n - 1));
    };  
  \end{lstlisting}
\end{adjustbox}
\end{multicols*}

\lstset{ 
  basicstyle=\ttfamily, 
  commentstyle=\color{blue}, 
  escapeinside={(*@}{@*)} 
}


\columnratio{.66,.33}
\begin{paracol}{2}
        \footnotesize







\begin{note}{}{}
    Les éléments de grammaire en \textbf{\textcolor{myr}{rouge}} 
    sont les élément manquants
\end{note}


\begin{adjustbox}{width=0.66\linewidth}
\begin{lstlisting}
(*@  \texttt{\textcolor{myb}{ \textit{/* Un programme correspond a une expression */  }}} @*)
<program> ::= <instr>


(*@  \texttt{\textcolor{myb}{ \textit{/* Une instruction est une expression suivis du terminal 
";" */  }}} @*)
<instr> ::= <expr> ";"
            (*@  \texttt{\textcolor{myb}{ \textit{/* Ou une instruction suivie d'une autre */  }}} @*)
            | <instr> <instr>

(*@  \texttt{\textcolor{myb}{ \textit{/* Une expression peut etre une fonction, une comparaison 
ou un appel de fonction */  }}} @*)
<expr> ::= "function" "(" <formals> ")" "{" <instr> "}"
            | <expr> "==" <expr>
            | <identifier> "(" <actuals> ")"
            (*@\texttt{\textcolor{red}{| <expr> + <expr>}  }   @*)
            (*@\texttt{\textcolor{red}{| <expr> - <expr>}  }   @*)
            (*@\texttt{\textcolor{red}{| <number>}  }   @*)
            (*@\texttt{\textcolor{red}{| "if" "(" <expr.> ")" "\{" <instr> "\}" <else> }}   @*)


(*@\texttt{\textcolor{red}{<else> :== "else" "\{" <instr.> "\}" }}   @*)
            (*@\texttt{\textcolor{red}{| "else if" "(" <expr.> ")" "\{" <instr.>"\}" <else> }}   @*)
            (*@\texttt{\textcolor{red}{| $\varepsilon $}}   @*)


            


(*@  \texttt{\textcolor{myb}{ \textit{/* Les parametres formels d'une fonction peuvent etre */  }}} @*)
<formals> ::= (*@ $\varepsilon$ @*) (*@ \texttt{ \textit{\textcolor{myb}{ vides}  }}@*)
             (*@ \texttt{ \textit{\textcolor{myb}{/* Ou un identifiant unique}  */}}@*) 
            | <identifier> { "," <identifier> }
             (*@ \texttt{ \textit{\textcolor{myb}{/* Ou une serie d'identifiants separes d'une virgule */}  }}@*)
            

(*@  \texttt{\textcolor{myb}{ \textit{/* Les parametre actuels d'un appel de fonction peuvent etre vide ou une expr. */  }}} @*)
<actuals> ::= (*@ $\varepsilon$ @*) | <expr>


(*@  \texttt{\textcolor{myb}{ \textit{/*  Permet a une expression d'etre un nombre */  }}} @*)

(*@  \texttt{\textcolor{myb}{ \textit{/* Un identifiant est une serie de lettres */  }}} @*)
(*@ \texttt{\textcolor{red}{<identifier> ::= <letter> \{letter\}}}  @*)

(*@  \texttt{\textcolor{myb}{ \textit{/* Les lettres acceptees sont celles de l'alphabet */  }}} @*)
(*@ \texttt{\textcolor{red}{<letter> ::= "a" - "z" | "A" - "Z"}}  @*)


(*@ \texttt{\textcolor{red}{<number> ::= <digit> \{digit\}}}  @*)


(*@ \texttt{\textcolor{red}{<digit> ::= |"0" | "1" | "2" | "3" | "4" | "5" | "6" | "7" | "8" | "9"}}  @*)


\end{lstlisting}






\end{adjustbox}


\switchcolumn

\end{paracol}












\end{document}
